\emph{(i)} Every element of $B_{\gamma}$ is a combination of basis blades $e_{S}$ (\claimref{std-algebra.blade-basis}), $S \subseteq \{0, \ldots, 4, \nu\}$, and for a generator $e_{i}$ of $A$ one has $e_{i}e_{S} = \pm e_{S}e_{i}$ blade by blade (\claimref{std-algebra.orthogonal-anticommute}): $e_{i}e_{S} = (-1)^{|S|}e_{S}e_{i}$ if $i \notin S$ and $(-1)^{|S|-1}e_{S}e_{i}$ if $i \in S$. Conjugation by $e_{i}$ therefore acts diagonally on the blade basis with eigenvalues $\pm 1$, and an element anticommutes with $e_{i}$ exactly when each of its blade components does. A blade anticommuting with all five $e_{i}$ either contains none of them, and then $|S|$ is odd, so $S = \{\nu\}$, or contains all five, and then $|S|$ is even, so $S = \{0, \ldots, 4, \nu\}$; a blade containing some but not all of them would need $|S|$ both odd and even. So $P_{\gamma} = \mathrm{span}\{\nu, e_{0}\cdots e_{4}\nu\} = \mathrm{span}\{\nu, \Omega\}$ with $\Omega = \omega\nu$. Since $\omega$ has odd grade and $\nu$ is orthogonal to its factors, $\omega\nu = -\nu\omega$, so $\nu\Omega + \Omega\nu = \nu\omega\nu + \omega\nu\nu = -\omega\gamma + \omega\gamma = 0$ and $\Omega^{2} = \omega\nu\omega\nu = -\omega^{2}\nu^{2} = -\gamma$. For $x = \alpha\nu + \beta\Omega$, $x^{2} = \gamma(\alpha^{2} - \beta^{2})$, a real form of signature $(1,1)$ for $\gamma = \pm 1$. Its null lines are $\alpha = \pm\beta$; $\nu - \Omega = \nu + \nu\omega = 2\nu\varepsilon_{+}$ and $\nu + \Omega = 2\nu\varepsilon_{-}$. Because $\omega$ anticommutes with $\nu$, $\nu\varepsilon_{\pm} = \varepsilon_{\mp}\nu$, which with $\varepsilon_{\pm}$ the diagonal units are the off-diagonal matrix units $E_{12} = \varepsilon_{+}\nu$, $E_{21} = \gamma\varepsilon_{-}\nu$ of Fact \claimref{rule2.imported-rule-2-matrix-units-forgotten} up to the factor $\gamma$; each has square $\nu\varepsilon_{\pm}\nu\varepsilon_{\pm} = \nu^{2}\varepsilon_{\mp}\varepsilon_{\pm} = 0$ and is nonzero. The check confirms the blade census in both signatures, the form on sample points and the matrix-unit relations $E_{12}E_{21} = \gamma\varepsilon_{+}$, $E_{21}E_{12} = \gamma\varepsilon_{-}$.

\emph{(ii)} Let $x \in P_{\gamma}$ with $x^{2} = \sigma \in \{\pm 1\}$. Then $e_{0}, \ldots, e_{4}, x$ pairwise anticommute with squares $1, -1, -1, -1, -1, \sigma$, so by the universal property there is a homomorphism from $\Cl(2,4)$ (if $\sigma = 1$) or $\Cl(1,5)$ (if $\sigma = -1$) onto the subalgebra they generate. Both of these Clifford algebras are simple of dimension $64$ (Proposition \claimref{rule2.rule-2-kaluza-klein-algebra}(iv)), so the homomorphism is injective and, $B_{\gamma}$ having dimension $64$, onto; the check also computes the span of the $64$ ordered products directly at five sample units of each sign in each $B_{\gamma}$ and finds rank $64$. The completion's volume element is $\omega x$: $\omega$ commutes with each $e_{i}$ and anticommutes with $x$, so $\omega x$ anticommutes with each $e_{i}$ (it lies in $P_{\gamma}$) and with $x$ (so it is orthogonal to $x$ for the polarised form $\tfrac12(xy + yx)$), and $(\omega x)^{2} = -\omega^{2}x^{2} = -\sigma$. For the isomorphism take $\gamma = -1$ and $x = \nu_{-}\omega$, which lies in $P_{-}$ with $x^{2} = -\nu_{-}^{2}\omega^{2} = 1$: the completion has signature $(2,4)$, so $e_{i} \mapsto e_{i}$, $\nu_{+} \mapsto \nu_{-}\omega$ is an isomorphism $\Cl(2,4) \to \Cl(1,5)$ fixing $A$ pointwise. It sends $\nu_{+}$ to $\nu_{-}\omega = -\Omega_{-}$ and $\Omega_{+} = \omega\nu_{+}$ to $\omega\nu_{-}\omega = -\nu_{-}$. The check verifies all $36$ anticommutation relations of the images, the rank $64$ of the image, and both images.

\emph{(iii)} If $\alpha$ is an automorphism with $\alpha(a) = a$ for $a \in A$, then $\alpha(x)$ anticommutes with every $e_{i}$ whenever $x$ does, so $\alpha(P_{\gamma}) = P_{\gamma}$, and $\alpha(x)^{2} = \alpha(x^{2}) = x^{2}$ because $x^{2}$ is a real scalar. So the sign of $x^{2}$ is invariant. The element $u = a\varepsilon_{+} + b\varepsilon_{-}$ with $ab \neq 0$ has inverse $a^{-1}\varepsilon_{+} + b^{-1}\varepsilon_{-}$ and commutes with $A$, since $\varepsilon_{\pm}$ are central in $A$. Using $\nu\varepsilon_{\pm} = \varepsilon_{\mp}\nu$, $u\,\nu\varepsilon_{+}\,u^{-1} = \nu(a\varepsilon_{-} + b\varepsilon_{+})\varepsilon_{+}(a^{-1}\varepsilon_{+} + b^{-1}\varepsilon_{-}) = (b/a)\,\nu\varepsilon_{+}$, and likewise $u\,\nu\varepsilon_{-}\,u^{-1} = (a/b)\,\nu\varepsilon_{-}$. In the null coordinates $x = p\,\nu\varepsilon_{+} + q\,\nu\varepsilon_{-}$ one has $x^{2} = pq(\nu\varepsilon_{+}\nu\varepsilon_{-} + \nu\varepsilon_{-}\nu\varepsilon_{+}) = pq\gamma$, and $\mathrm{Ad}_{u}$ sends $(p, q)$ to $(tp, q/t)$ with $t = b/a$ ranging over all nonzero reals; this group acts transitively on each hyperbola $pq = c/\gamma$, both branches included since $t$ may be negative. Every unit vector completing $A$'s generators lies in $P_{\gamma}$ (it must anticommute with them), so the completions fall into exactly two classes, one per sign, and no automorphism fixing $A$ joins them. The check verifies the conjugation formula for three pairs $(a,b)$ of mixed signs in each $B_{\gamma}$ and exhibits the $u$ that carries one point of a level set to another on the opposite branch.

\emph{(iv)} A property of the pair $(B, A)$ invariant under automorphisms takes the same value on two completions in one class by (iii); within one algebra both classes are present by (ii); so the only invariant datum that separates the two repairs is the class, the sign of $x^{2}$. By (ii) the completion's volume element squares to $-x^{2}$. The even part of a Clifford algebra on six generators has centre $\mathrm{span}\{1, \omega x\}$ (Proposition \claimref{synthesis.plane-commutant-even-ground}(ii), where the centre of the even algebra is $\mathrm{span}\{1, \omega\}$ in even dimension), which is $\C$ exactly when the volume element squares to $-1$ and $\R \oplus \R$ when it squares to $+1$; the check verifies that the volume element commutes with all fifteen bivectors of each completion. For $\Cl(2,4)$ this volume element is the complex unit of $\Cl^{0}(2,4) \cong \Mat(4,\C)$, as Proposition \claimref{odd-extensions.twistor-algebra}(i) states. The three statements, $x^{2} = 1$, the volume element squares to $-1$, and the even part has complex centre, are therefore one condition, and it is the residue $6$ against $4$ of Proposition \claimref{rule2.rule-2-kaluza-klein-algebra}(vi) read modulo four. Nothing in this argument decides the condition.
